\section{Appendix: Bias, covariance, profiling, and attainment}

The calibration in \cref{obj:thm:whole-null-calibration-resolved} rests on three population quantities: the mean of the original-record score, the covariance of its affine coefficients, and the fluctuation of the cross-block inner product uniformly over candidate constants. This appendix organizes those quantities around the public construction in \cref{obj:def:explicit-score-ledger}. The resulting bounds connect the finite score calculation to the explicit separation guarantee.

\subsection{Calibration and population decomposition}

The deterministic calibration handle collects the score construction, its affine coefficients, the public bounds, and the profiling operation. Keeping these components together makes their dependence on the sample size and public exponent tuple explicit.

\begin{algorithmv}{def:calibration-handle}[Deterministic calibration handle]
The calibration handle, indexed by a sample size \(n\in\mathbb N\) and a public exponent tuple \(v=(p,\alpha,\beta,\gamma)\in\mathbb R^4\), is the ordered collection of the following seven components. Its branch, ranks, cutoffs, ledgers, and public grids are the deterministic selections from \(n\) and \(v\) in \cref{obj:def:explicit-score-ledger}.

\begin{enumerate}
\item The tuning collection
\[
\bigl(M,K,T_0,(T_j)_{j=1}^{L}\bigr),
\qquad M=J_{\mathrm c},\quad K=J_{\mathrm f},
\]
with the number of increments \(L\) specified in \cref{obj:def:explicit-score-ledger}.

\item The block-score family, taking a block index, a real candidate constant \(c\), and a dataset of \(n\) original records to a vector in \(\mathbb R^M\). For \(n\ge4\), it is
\[
Z_b(c)=
\begin{cases}
H_{b,T_0}(c)-U_{b,\Pi_K,T_0}(c),
&\text{single-correction branch},\\[1mm]
H_{b,T_0}(c)-U_{b,\Pi_M,T_0}(c)
-\displaystyle\sum_{j=1}^{L}U_{b,D_j,T_j}(c),
&\text{multiresolution branch},
\end{cases}
\]
where the block and kernel conventions are those of
\cref{obj:def:blocks,obj:def:score-family,obj:def:explicit-score-ledger}.
For \(n<4\), \(Z_b(c)=0\).

\item The two affine-coefficient maps, indexed by the block and a binary coefficient selector, with values
\[
Z_{b,0}=Z_b(0),
\qquad
Z_{b,A}=Z_b(0)-Z_b(1).
\]
Thus, for every real \(c\) and every dataset,
\[
Z_b(c)=Z_{b,0}-cZ_{b,A}.
\]

\item The deterministic ledger triple
\[
\bigl(\mathfrak b_{n,v},\Lambda_{n,v},h_{n,v}\bigr),
\]
whose bias, covariance, and rejection-cutoff selections are exactly those of
\cref{obj:def:explicit-score-ledger}.

\item The dataset-to-real profiling map
\[
\inf_{c\in[-1/2,1/2]}W(c),
\qquad
W(c)=\langle Z_1(c),Z_2(c)\rangle,
\]
using the two-block indexing convention of
\cref{obj:def:explicit-score-ledger}.

\item The original-record rejection rule specified in
\cref{obj:def:explicit-score-ledger}, namely, for \(n\ge4\),
\[
\mathbf 1\!\left\{
\inf_{c\in[-1/2,1/2]}W(c)>h_{n,v}
\right\},
\]
and the zero rejection rule for \(n<4\).

\item The ordered triple of finite public rank, outcome-cutoff, and rejection-cutoff grids selected in \cref{obj:def:explicit-score-ledger}. The first two grids depend on \(n\), and the third depends on \(n\) and \(v\).
\end{enumerate}
\end{algorithmv}

The affine coefficient \(Z_{b,A}\) records the score's response to the candidate constant, while \(Z_{b,0}\) records its outcome contribution. Their joint population behavior is therefore sufficient to describe the entire profiled statistic. The mean handle below retains both coefficients, their singleton and canonical pair projections, and all four combinations of single-record and pair contributions.

\begin{algorithmv}{def:mean-handle}[Original score mean handle]
The original score mean handle is the ordered pair of multiresolution and single-correction handles associated with the two score members of \cref{obj:def:calibration-handle}. Its parameters are an original-record law \(P\), ranks \(J,K,L\in\mathbb N\), a main cutoff \(T_{\mathrm{main}}\in\mathbb R\), and increment cutoffs \(T_j^{\mathrm{inc}}\in\mathbb R\) indexed by \(1\le j\le L\). Each branch retains, in order, a mean map indexed by \(c\in\mathbb R\), coefficient means, singleton projections, canonical pair projections, four population contractions for each coefficient pair, and a sample statistic indexed by \(n\) and \(c\).

The multiresolution branch uses coarse rank \(J_{\mathrm c}=J\) and fine rank \(J_{\mathrm f}=2^LJ\); the single-correction branch uses
\[
J_{\mathrm c}=J,\qquad J_{\mathrm f}=K.
\]
Both branches use \(T_{\mathrm{main}}\); the multiresolution branch also uses the increment cutoffs \(T_j^{\mathrm{inc}}\), \(1\le j\le L\). Define the projected propensity and propensity weight by
\[
e_K=\Pi_{J_{\mathrm f}}e_P,
\qquad
w_K(x)=e_P(x)\{1-e_K(x)\}.
\]
The retained mean map is
\[
\int_0^1 F_{J_{\mathrm c}}(x)
\left[
w_K(x)\{\tau_P(x)-c\}
+\{e_P(x)-e_K(x)\}m_{0,P}(x)
\right]\,dx
\]
plus the exact multiresolution or single-correction clipping-tail remainder, respectively, specified in \cref{obj:lem:original-score-mean-ledger}.

This retained map equals the population expectation of the corresponding corrected block score for every \(c\in\mathbb R\) and either evaluation block under the following conditions:
\begin{itemize}
\item \textbf{(Model.)} The public tuple \(v\) is admissible and \(P\in\mathcal M_v\), as specified in \cref{obj:def:model}.
\item \textbf{(Sample and coarse rank.)} \(n\ge4\) and \(J>0\).
\item \textbf{(Main cutoff.)} \(T_{\mathrm{main}}\ge1\).
\item \textbf{(Multiresolution cutoffs.)} In the multiresolution branch, \(T_j^{\mathrm{inc}}\ge1\) for every integer \(1\le j\le L\).
\item \textbf{(Single-correction refinement.)} In the single-correction branch, \(K>0\) and \(J\mid K\).
\end{itemize}

For each affine coefficient \(r\in\{0,1\}\), let \(h_r\) and \(g_r\) be the selected branch's raw single-record and symmetric pair kernels from \cref{obj:def:calibration-handle}. All expectations below use independent original records with law \(P\). Define the centered pair-kernel singleton projection by
\[
g_{r,1}(o)
=
\mathbb E[g_r(o,O')]
-\mathbb E[g_r(O,O')].
\]
The retained coefficient means, singleton projections, and canonical pair projections are, respectively,
\[
\mathbb E[h_r(O)]+\mathbb E[g_r(O,O')],
\]
\[
h_r(o)-\mathbb E[h_r(O)]+2g_{r,1}(o),
\]
\[
g_r(o,z)-\mathbb E[g_r(O,O')]
-g_{r,1}(o)-g_{r,1}(z).
\]
For every \(r,r'\in\{0,1\}\), the retained ordered quadruple of population contractions is defined by
\[
h_r\!\times h_{r'}
=
\mathbb E\langle h_r(O_1),h_{r'}(O_2)\rangle,
\]
\[
h_r\!\times g_{r'}
=
\mathbb E\langle h_r(O_1),g_{r'}(O_2,O_3)\rangle,
\]
\[
g_r\!\times h_{r'}
=
\mathbb E\langle g_r(O_1,O_2),h_{r'}(O_3)\rangle,
\]
\[
g_r\!\times g_{r'}
=
\mathbb E\langle g_r(O_1,O_2),g_{r'}(O_3,O_4)\rangle,
\]
where \(O_1,O_2,O_3,O_4\) are independent with law \(P\).

For each candidate constant \(c\in\mathbb R\), set
\[
h_c=h_0-ch_1,
\qquad
g_c=g_0-cg_1.
\]
For \(n\ge4\), use the block size \(s_n=\lfloor n/2\rfloor\) and the two deterministic blocks \(\mathcal I_1,\mathcal I_2\) of \cref{obj:def:blocks}. Define
\[
Z_b(c)=
\frac1{s_n}\sum_{i\in\mathcal I_b}h_c(O_i)
+
\frac1{s_n(s_n-1)}
\sum_{\substack{i,j\in\mathcal I_b\\i\ne j}}
g_c(O_i,O_j),
\qquad b\in\{1,2\}.
\]
The retained sample output is the single function \(c\mapsto W(c)\), given by
\[
\begin{aligned}
W(c)
&=\langle Z_1(c),Z_2(c)\rangle\\
&=\frac1{s_n^2}
\sum_{\substack{i\in\mathcal I_1\\j\in\mathcal I_2}}
\langle h_c(O_i),h_c(O_j)\rangle\\
&\quad+\frac1{s_n^2(s_n-1)}
\sum_{\substack{i\in\mathcal I_1\\j,k\in\mathcal I_2,\ j\ne k}}
\langle h_c(O_i),g_c(O_j,O_k)\rangle\\
&\quad+\frac1{s_n^2(s_n-1)}
\sum_{\substack{i,j\in\mathcal I_1,\ i\ne j\\k\in\mathcal I_2}}
\langle g_c(O_i,O_j),h_c(O_k)\rangle\\
&\quad+\frac1{s_n^2(s_n-1)^2}
\sum_{\substack{i,j\in\mathcal I_1,\ i\ne j\\k,l\in\mathcal I_2,\ k\ne l}}
\langle g_c(O_i,O_j),g_c(O_k,O_l)\rangle.
\end{aligned}
\]
For \(n<4\), its defining value is
\[
W(c)=0.
\]
The two mixed affine-coefficient contributions are retained through their sum in \(W(c)\).
\end{algorithmv}

The conventions in \cref{obj:def:mean-handle,obj:def:calibration-handle} use the same one-based blocks \(\mathcal I_1,\mathcal I_2\). In the mean calculation, \(T_{\mathrm{main}}\) denotes the procedure's main cutoff \(T_0\), while \(T_j^{\mathrm{inc}}\) denotes its increment cutoff \(T_j\) for \(1\le j\le L\). Thus the two cutoff families and the two block scores have distinct names without a positional shift. At the public selections, the coarse rank is \(J=M\) and the multiresolution fine rank is \(2^LJ=K\).

For clarity about the symmetric pair convention, write \(o=(x,a,y)\) and \(z=(x',a',y')\). For a correction function \(G\) and cutoff \(T\), define its signed, symmetrized contribution by
\[
k_{G,T,c}(o,z)
:=-\frac12\left[
F_M(x)aG(x,x')\{\ell_T(y')-ca'\}
+
F_M(x')a'G(x',x)\{\ell_T(y)-ca\}
\right].
\]
The single-record contribution is
\[
h_c(o):=F_M(x)a\{\ell_{T_0}(y)-c\}.
\]
In the single-correction branch, \(g_c:=k_{\Pi_K,T_0,c}\). In the multiresolution branch, using the increment indices of \cref{obj:def:explicit-score-ledger},
\[
g_c:=k_{\Pi_M,T_0,c}
+\sum_{j=1}^{L}k_{D_j,T_j,c}.
\]
The affine coefficients in \cref{obj:def:mean-handle} are thus \(h_0=h_c|_{c=0}\), \(h_1=h_c|_{c=0}-h_c|_{c=1}\), and the corresponding two coefficients of \(g_c\). Symmetrization preserves the ordered-pair average because that average contains both orientations of every distinct pair.

The four contractions in \cref{obj:def:mean-handle} retain the single--single, single--pair, pair--single, and pair--pair contributions separately. Their independent-record arguments match the four finite sums defining \(W(c)\). Both orientations of the mixed contribution are retained, including their dependence on the two affine coefficients.

\subsection{Original-mean bias}

The population calculation separates a weighted effect contrast, a baseline approximation term, and clipping tails. To specify the latter explicitly, let \(\lambda\) denote Lebesgue probability measure on \([0,1]\), and define the conditional clipping discrepancy and its unconditional conditional mean by
\[
\delta_{a,T}(x)
:=\mathbb E_P[\ell_T(Y)-Y\mid X=x,A=a],
\qquad
r_T(x):=e_P(x)\delta_{1,T}(x)
+\{1-e_P(x)\}\delta_{0,T}(x).
\]
These quantities are defined up to \(\lambda\)-almost-everywhere equality; their integrated contributions use that convention. Under the mean handle's conditions, the clipping-tail remainders for the two branches are
\[
\begin{aligned}
\mathcal T_{\mathrm{single}}
&:=\int_0^1 F_M(x)
\left[e_P(x)\delta_{1,T_0}(x)
-e_K(x)r_{T_0}(x)\right]\,d\lambda(x),\\
\mathcal T_{\mathrm{multi}}
&:=\int_0^1 F_M(x)
\left[
e_P(x)\delta_{1,T_0}(x)
-(\Pi_Me_P)(x)r_{T_0}(x)
-\sum_{j=1}^{L}(D_je_P)(x)r_{T_j}(x)
\right]\,d\lambda(x).
\end{aligned}
\]
Here \((D_je_P)(x):=\int_0^1D_j(x,z)e_P(z)\,d\lambda(z)\), and the increment indices follow \cref{obj:def:explicit-score-ledger}. Each remainder is independent of the candidate constant.

The retained mean in \cref{obj:def:mean-handle} therefore identifies the baseline approximation contribution through \((e_P-e_K)m_{0,P}\), while the clipping remainder records the outcome-tail contribution at each selected resolution. The public bound \(B\) in \cref{obj:def:explicit-score-ledger} aggregates these terms. The next statement controls that bound uniformly over the candidate constants and relates the score mean to heterogeneity.

\begin{lemmav}{lem:original-score-mean-ledger}[Original score mean bounds]
Let \(v=(p,\alpha,\beta,\gamma)\) be an admissible tuple and \(P\in\mathcal M_v\), as specified in \cref{obj:def:model}, and let \(n\ge4\) be an integer. Use the explicit score ledger's coarse rank \(M\), fine rank \(K\), corrected block scores \(Z_b(c)\), and deterministic bias bound \(B\).

Let \(\lambda\) denote Lebesgue probability measure on \([0,1]\), and define the projected propensity and its associated weight by
\[
e_K(x)=(\Pi_K e_P)(x),
\qquad
w_K(x)=e_P(x)\bigl(1-e_K(x)\bigr).
\]
Define the population heterogeneity distance by
\[
d(P)=
\left(
\int_0^1
\left[\tau_P(x)-\int_0^1\tau_P(z)\,d\lambda(z)\right]^2
\,d\lambda(x)
\right)^{1/2}.
\]

For every block index \(b\) and every \(c\in\mathbb R\), \(Z_b(c)\) is integrable under the original independent-record law \(P^{\otimes n}\). Its population mean is
\[
\theta_P(c)
=
\mathbb E_{P^{\otimes n}}\!\left[Z_b(c)\right].
\]
For every block index \(b\) and every \(c\in\mathbb R\) with \(|c|\le\tfrac12\),
\[
\left\|
\theta_P(c)
-
\int_0^1
w_K(x)\bigl(\tau_P(x)-c\bigr)F_M(x)\,d\lambda(x)
\right\|_2
\le B,
\qquad
\|\theta_P(c)\|_2
\ge
\frac{3}{16}d(P)-16M^{-\gamma}-B.
\]
Furthermore, if \(P\in H_0(v)\) as defined in \cref{obj:def:null}, then for every block index \(b\) and every \(c_0\in\mathbb R\) satisfying \(\tau_P(x)=c_0\) for all \(x\in[0,1]\),
\[
\|\theta_P(c_0)\|_2\le B.
\]
\end{lemmav}

The mean bounds supply the population distinction used by the test. Under a constant effect, evaluating at that constant gives a score mean of norm at most \(B\). Under heterogeneity, every admissible candidate constant has a score mean bounded below by a multiple of \(d(P)\), after accounting for coarse approximation and original-mean bias. Uniformity in the candidate constant is what makes these bounds compatible with profiling.

\subsection{Singleton, pair, and cross-level covariance bounds}

The singleton projection in \cref{obj:def:mean-handle} combines the centered single-record contribution with twice the pair-kernel singleton projection. Its canonical pair projection retains the component centered with respect to either record. These two components have separate public bounds, \(L_1\) and \(L_2\), in \cref{obj:def:explicit-score-ledger}.

For the multiresolution construction, both bounds aggregate the main correction and every retained increment. In particular, \(L_1\) contains the sum of the increment-specific smoothness and clipping contributions, while \(L_2\) contains the sum of their rank-weighted second-moment bounds. The squared sums in
\[
\Lambda_{n,v}
=\frac{L_1^2}{s}
+\frac{2L_2^2}{s(s-1)}
\]
retain the cross-level contributions within each component. The covariance statement below controls both affine coefficients and their covariance in arbitrary, potentially different, directions.

\begin{lemmav}{lem:original-score-covariance-ledger}[Original score covariance bounds]
Let \(v=(p,\alpha,\beta,\gamma)\) be an admissible tuple of moment and smoothness exponents, and let \(P\in\mathcal M_v\), as specified in \cref{obj:def:model}. Let \(n\ge 2\) be an integer. Write the affine block score as
\[
Z_b(c)=Z_{b,0}-cZ_{b,A}.
\]
Under the independent-record sample law \(P^{\otimes n}\), for every evaluation block \(b\) and every \(u\in\mathbb R^M\), both \(u^T Z_{b,0}\) and \(u^T Z_{b,A}\) are square integrable and satisfy
\[
\operatorname{Var}_{P^{\otimes n}}(u^T Z_{b,0})
\le \Lambda_{n,v}\|u\|_2^2,
\qquad
\operatorname{Var}_{P^{\otimes n}}(u^T Z_{b,A})
\le \Lambda_{n,v}\|u\|_2^2.
\]
Moreover, for every evaluation block \(b\) and every \(u,z\in\mathbb R^M\),
\[
\left|
\operatorname{Cov}_{P^{\otimes n}}
\bigl(u^T Z_{b,0},z^T Z_{b,A}\bigr)
\right|
\le \Lambda_{n,v}\|u\|_2\|z\|_2.
\]
\end{lemmav}

The directional form gives a common covariance bound for the outcome coefficient and the constant-effect coefficient. The mixed bound retains their dependence within a block, including the dependence contributed by the different correction levels. Independence between the two evaluation blocks then provides the setting for the cross-block profiling statement.

\subsection{Simultaneous affine profiling and calibration}

The profiling multiplier \(C=128\) and the population mean \(\theta_P(c)\) were introduced with \cref{obj:thm:whole-null-calibration-resolved}. The following bound applies to the entire interval of candidate constants on one measurable event.

\begin{lemmav}{lem:uniform-affine-profile-bound}[Uniform affine profile bound]
Let \(v=(p,\alpha,\beta,\gamma)\) be an admissible public tuple as in \cref{obj:def:model}, and let \(n\ge 2\) be an integer. For every \(P\in\mathcal M_v\), consider the explicit affine score construction, with block scores \(Z_1(c)\) and \(Z_2(c)\), coarse histogram rank \(J_{\mathrm c}\), and public covariance ledger \(\Lambda_{n,v}\). Define the score inner product and its population target by
\[
W(c)=\langle Z_1(c),Z_2(c)\rangle,
\qquad
\theta_P(c)=\mathbb E_{P^{\otimes n}}[Z_1(c)].
\]
The event on which, simultaneously for every \(c\in[-1/2,1/2]\),
\[
\bigl|W(c)-\|\theta_P(c)\|_2^2\bigr|
\le
128\left\{
\sqrt{\Lambda_{n,v}}\,\|\theta_P(c)\|_2
+\sqrt{J_{\mathrm c}}\,\Lambda_{n,v}
\right\}
\]
is measurable and has probability at least \(0.9925\) under \(P^{\otimes n}\). Moreover, if \(\Lambda_{n,v}=0\), then \(P^{\otimes n}\)-almost surely, simultaneously for every \(c\in[-1/2,1/2]\),
\[
W(c)=\|\theta_P(c)\|_2^2.
\]
\end{lemmav}

Simultaneous control permits minimization over the candidate constant within the same event. Its fluctuation bound separates a term proportional to the population score norm from a term involving the coarse dimension and covariance ledger. When the covariance ledger vanishes, the simultaneous identity gives the population target exactly.

The public computation retains this affine structure explicitly. Using the one-based block convention of \cref{obj:def:calibration-handle}, set \(Z_{b,1}:=Z_{b,A}\) and
\[
w_{rr'}:=\langle Z_{1,r},Z_{2,r'}\rangle,
\qquad r,r'\in\{0,1\}.
\]
Then the computed quadratic is
\[
W(c)=w_{00}-c(w_{01}+w_{10})+c^2w_{11}.
\]
Its minimum over \([-1/2,1/2]\) is obtained by comparing the two endpoints and, when \(w_{11}>0\) and the vertex lies in the interval, the value at
\[
c_{\mathrm{vert}}:=\frac{w_{01}+w_{10}}{2w_{11}}.
\]
This formula preserves both mixed affine-coefficient contributions through their sum. The same calculation applies under the zero-based convention after translating the block indices.

Together, \cref{obj:lem:original-score-mean-ledger,obj:lem:original-score-covariance-ledger,obj:lem:uniform-affine-profile-bound} supply the population and fluctuation bounds used in \cref{obj:thm:whole-null-calibration-resolved}. The cutoff in \cref{obj:def:explicit-score-ledger} rounds
\[
2B^2+2^{16}\sqrt M\,\Lambda_{n,v}
\]
upward to a dyadic value. The calibration theorem gives size at most \(0.0075\) throughout the constant-effect null and type-II error at most \(0.0075\) at its calibrated separation, for \(n\ge4\) when that separation is below the model distance supremum. The zero rejection rule supplies the stated size guarantee at \(n=2,3\).

The calibrated separation \(R_{\mathrm{att}}(n,v)\) in \cref{obj:thm:whole-null-calibration-resolved} groups coarse approximation, the bias ledger, and the covariance contribution. The explicit attainment statement below expresses their combined requirement as a power of the sample size with a fully specified multiplier.

\subsection{Explicit attainment constants and separation}

Fix an admissible tuple \(v=(p,\alpha,\beta,\gamma)\) from \cref{obj:def:model}, and use the exponents \(q,S,t,E_0,E_4,E\) of \cref{obj:def:explicit-score-ledger}. Before stating attainment, define the minimum smoothness \(s_0\), the clipping allocation exponent \(\lambda_0\), the increment summability exponent \(\eta_0\), and the geometric-series factor \(A(x)\) by
\[
s_0:=\min\{\alpha,\beta,\gamma\},
\qquad
\lambda_0:=\frac{(p-1)^2}{4p},
\qquad
\eta_0:=\frac{(p-1)(p+6)}{4p},
\qquad
A(x):=(1-2^{-x})^{-1}\quad(x>0).
\]
Here \(A(x)\) denotes a scalar series factor; its argument distinguishes it from the treatment coordinate. Define the grouped bias, summation, singleton, and canonical-pair constants by
\[
\begin{aligned}
B_*&:=420+800\{A(\alpha)+A(\lambda_0)\},\\
D_*&:=1+\{\exp(1)S\log 2\}^{-1},\\
L_*&:=16+880A(s_0)+160D_*+960A(\alpha),\\
G_*&:=\sqrt{64}\{2^{t/2}A(1/2)+A(\eta_0/2)\}.
\end{aligned}
\]
The aggregate covariance constant \(A_*\) and attainment multiplier \(C^{\mathrm{att}}_v\) are
\[
A_*:=\sqrt{2}\{24576+64L_*^2+16384+256G_*^2\},
\qquad
C^{\mathrm{att}}_v:=16\{16+5B_*+1024\sqrt{A_*}\}.
\]
These definitions resolve every constant entering the public attainment multiplier.

Using the heterogeneity distance and model supremum defined in \cref{obj:thm:whole-null-calibration-resolved}, define the fixed witness distance \(d_0\) and public threshold \(N^{\mathrm{att}}_v\) by
\[
d_0:=\frac{1}{8\sqrt{2}},
\qquad
N^{\mathrm{att}}_v
:=\left\lceil\max\left\{4,
\left(\frac{2C^{\mathrm{att}}_v}{d_0}\right)^{1/E}\right\}\right\rceil.
\]
For the rejection expectation below, \(\mathcal D_n\) denotes \(n\) independent original records with law \(P\), and \(U\) denotes an independent uniform variable on \([0,1]\), as in \cref{obj:def:original-record-experiment}. Set
\[
\mathsf R_n(P,\phi):=\mathbb E[\phi(\mathcal D_n,U)].
\]

For an admissible bounded-model tuple \(w=(\alpha,\beta,\gamma)\) from \cref{obj:def:bounded-model}, define the matched exponent, multiplier, and distance supremum by
\[
E_{\mathrm b}
:=\min\left\{\frac{2\gamma}{4\gamma+1},
\frac{2(\alpha+\beta)}
{1+2(\alpha+\beta)+(\alpha+\beta)/(2\gamma)}\right\},
\qquad
C^{\mathrm{att}}_{(2,w)}
:=C^{\mathrm{att}}_{(2,\alpha,\beta,\gamma)},
\qquad
D_w^{\mathrm b}:=\sup_{P\in\mathcal M_w^{\mathrm b}}d(P).
\]
The following result gives attainment for the original model and its matched bounded submodel, together with public tuning bounds and uniform control of the constants on compact sets of admissible tuples.

\begin{theoremv}{thm:original-record-attainable-rate}[Original-record attainable separation rate]
For every admissible public tuple \(v=(p,\alpha,\beta,\gamma)\) of \cref{obj:def:model}, use the explicit rule \(\phi^*_{n,v}\) and its tuning quantities from \cref{obj:def:explicit-score-ledger,obj:def:sharp-frontier-scales}. Write
\[
q=\frac{p-1}{p},\qquad S=\alpha+\beta,\qquad
t=\frac{2-p}{p-1},\qquad
E_0=\frac{2\gamma q}{2\gamma+q},\qquad
E_4=\frac{2S}{1+2\alpha+\beta/q+S/(2\gamma)},\qquad
E=\min\{E_0,E_4\}.
\]
Define the minimum smoothness \(s_0\), increment summability exponent \(\eta_0\), geometric-series factor \(A(x)\), and attainment constants by
\[
\begin{aligned}
s_0&=\min\{\alpha,\beta,\gamma\},&
\lambda_0&=\frac{(p-1)^2}{4p},&
\eta_0&=\frac{(p-1)(p+6)}{4p},\\
A(x)&=(1-2^{-x})^{-1}\quad(x>0),\\
B_*&=420+800\{A(\alpha)+A(\lambda_0)\},\\
D_*&=1+\{\exp(1)S\log 2\}^{-1},\\
L_*&=16+880A(s_0)+160D_*+960A(\alpha),\\
G_*&=\sqrt{64}\{2^{t/2}A(1/2)+A(\eta_0/2)\},\\
A_*&=\sqrt{2}\{24576+64L_*^2+16384+256G_*^2\},\\
C^{\mathrm{att}}_v&=16\{16+5B_*+1024\sqrt{A_*}\}.
\end{aligned}
\]
Define the fixed witness distance \(d_0\) and public attainment threshold \(N^{\mathrm{att}}_v\) by
\[
d_0=\frac{1}{8\sqrt{2}},
\qquad
N^{\mathrm{att}}_v
=\left\lceil\max\left\{4,
\left(\frac{2C^{\mathrm{att}}_v}{d_0}\right)^{1/E}\right\}\right\rceil.
\]
Let \(\lambda\) be Lebesgue probability measure on \([0,1]\). Define the heterogeneity distance \(d(P)\) and its model supremum \(D_v\) by
\[
d(P)=
\left\{\int_{[0,1]}
\left(\tau_P(x)-\int_{[0,1]}\tau_P(z)\,d\lambda(z)\right)^2
\,d\lambda(x)\right\}^{1/2},
\qquad
D_v=\sup_{P\in\mathcal M_v}d(P).
\]
Let \(\mathcal D_n\) be a sample of \(n\) independent original records with law \(P\), and let \(U\) be an independent uniform randomization variable on \([0,1]\). For a test \(\phi\), define its rejection expectation by
\[
\mathsf R_n(P,\phi)=\mathbb E[\phi(\mathcal D_n,U)].
\]

Then \(C^{\mathrm{att}}_v>0\), and the following guarantees hold:
\begin{itemize}
\item \textbf{(Size and separation.)}
For every integer \(n\ge2\),
\[
\mathsf R_n(P,\phi^*_{n,v})\le0.0075
\qquad\text{for every }P\in H_0(v),
\]
where the null is given in \cref{obj:def:null}. Whenever
\(C^{\mathrm{att}}_v n^{-E}<D_v\),
\[
1-\mathsf R_n(P,\phi^*_{n,v})\le0.0075
\qquad
\text{for every }P\in\mathcal M_v
\text{ with }d(P)\ge C^{\mathrm{att}}_v n^{-E}.
\]
The capped critical radius of \cref{obj:def:critical-radius} satisfies
\[
r_n^*(v)\le C^{\mathrm{att}}_v n^{-E}.
\]

\item \textbf{(Witness threshold.)}
For every integer \(n\ge N^{\mathrm{att}}_v\),
\[
C^{\mathrm{att}}_v n^{-E}<d_0.
\]

\item \textbf{(Tuning bounds.)}
For every integer \(n\ge4\), the coarse rank \(M\), fine rank \(K\), block size \(s\), main clipping cutoff \(T_0\), and increment cutoffs \(T_j\) from \cref{obj:def:explicit-score-ledger,obj:def:blocks} satisfy
\[
M\le2s,\qquad K\le2s^2,\qquad
1\le T_0\le2s,\qquad
1\le T_j\le T_0\quad(1\le j\le L).
\]

\item \textbf{(Compact uniformity.)}
For every compact set \(K\) of public tuples in the product topology whose members are all admissible under \cref{obj:def:model}, there exist real constants \(C,N,e\), with \(e>0\), such that, for every \(v\in K\),
\[
C^{\mathrm{att}}_v\le C,\qquad
N^{\mathrm{att}}_v\le N,\qquad e\le E(v).
\]
Here \(E(v)\) denotes the displayed minimum exponent evaluated at \(v\).
\end{itemize}

For every admissible smoothness tuple \(w=(\alpha,\beta,\gamma)\) of \cref{obj:def:bounded-model}, define the bounded separation exponent \(E_{\mathrm b}\), matched attainment multiplier \(C^{\mathrm{att}}_{(2,w)}\), and bounded-model distance supremum \(D_w^{\mathrm b}\) by
\[
E_{\mathrm b}
=\min\left\{\frac{2\gamma}{4\gamma+1},
\frac{2(\alpha+\beta)}
{1+2(\alpha+\beta)+(\alpha+\beta)/(2\gamma)}\right\},
\qquad
C^{\mathrm{att}}_{(2,w)}
=C^{\mathrm{att}}_{(2,\alpha,\beta,\gamma)},
\qquad
D_w^{\mathrm b}=\sup_{P\in\mathcal M^{\mathrm b}_w}d(P).
\]
The original-model exponent evaluated at \((2,\alpha,\beta,\gamma)\) equals \(E_{\mathrm b}\). The bounded-model rule \(\phi^{*,\mathrm b}_{n,w}\) from \cref{obj:def:sharp-frontier-scales}, obtained by evaluating the same ledger rule at \((2,\alpha,\beta,\gamma)\), satisfies, for every integer \(n\ge2\),
\[
\mathsf R_n(P,\phi^{*,\mathrm b}_{n,w})\le0.0075
\qquad\text{for every }P\in H_0^{\mathrm b}(w),
\]
with the bounded null as in \cref{obj:def:bounded-null}. Whenever
\(C^{\mathrm{att}}_{(2,w)}n^{-E_{\mathrm b}}<D_w^{\mathrm b}\),
\[
1-\mathsf R_n(P,\phi^{*,\mathrm b}_{n,w})\le0.0075
\qquad
\text{for every }P\in\mathcal M^{\mathrm b}_w
\text{ with }d(P)\ge C^{\mathrm{att}}_{(2,w)}n^{-E_{\mathrm b}}.
\]
The bounded capped critical radius of \cref{obj:def:bounded-radius} satisfies
\[
r_n^{*,\mathrm b}(w)
\le C^{\mathrm{att}}_{(2,w)}n^{-E_{\mathrm b}}.
\]
\end{theoremv}

The attainment guarantee combines the mean separation in \cref{obj:lem:original-score-mean-ledger} with the simultaneous fluctuation bound in \cref{obj:lem:uniform-affine-profile-bound}. Its exponent is the minimum of the two public separation exponents, while its multiplier accounts explicitly for projection, clipping, and covariance contributions. Evaluating the same construction at moment order two gives the bounded-model exponent and multiplier. The power clauses apply when the stated separation is below the corresponding model distance supremum; the capped-radius bounds apply with the capping conventions in \cref{obj:def:critical-radius,obj:def:bounded-radius}.

The tuning bounds place every selected rank and outcome cutoff within the finite public grids of \cref{obj:def:explicit-score-ledger}. Compact uniformity supplies common bounds for the attainment constants and thresholds over compact sets of admissible tuples. Thus the explicit procedure links a finite original-record calculation to a separation guarantee whose dependence on the public moment and smoothness parameters is fully specified.
